\section{2012级工科数学分析下 B卷}

\subsection{资料来源}
\begin{itemize}
  \item 试题：2012级计算机工科数学分析下期末考试题 B 卷，DOC 正文已抽取；公式对象已按 Word 公式图片逐项恢复。
  \item 公式图片审计：\texttt{tmp/past-exam-equations/2012-B/manifest.csv}。
\end{itemize}

\subsection{试题}
诚信应考，考试作弊将带来严重后果！

华南理工大学期末考试

《工科数学分析》2012--2013 第二学期期末考试试卷（B）

注意事项：
\begin{enumerate}
  \item 考前请将密封线内填写清楚；
  \item 所有答案请直接答在试卷上（或答题纸上）；
  \item 考试形式：闭卷；
  \item 本试卷共 5 个大题，满分 100 分，考试时间 120 分钟。
\end{enumerate}

\subsubsection*{一、单项选择题（每小题 3 分，共 15 分）}
\begin{enumerate}
  \item 初值问题
  \[
  y''+4y'+4y=0,\qquad y(0)=1,\quad y'(0)=3
  \]
  的解为（\quad）。

  A. $y=(C_1+C_2x)e^{-2x}$；\quad B. $y=(1+5x)e^{-2x}$；\quad
  C. $y=(1+3x)e^{-2x}$；\quad D. $y=e^{-2x}+5e^{2x}$。

  \item 对于二元函数 $z=f(x,y)$，下列说法正确的是（\quad）。

  A. 若函数 $z=f(x,y)$ 在 $(x_0,y_0)$ 连续，则函数在该点一定存在偏导数；

  B. 若函数 $z=f(x,y)$ 在 $(x_0,y_0)$ 存在偏导数，则函数在该点一定连续；

  C. 若函数 $z=f(x,y)$ 在 $(x_0,y_0)$ 的偏导数连续，则函数在该点一定可微；

  D. 若函数 $z=f(x,y)$ 在 $(x_0,y_0)$ 可微，则函数在该点的偏导数一定连续。

  \item 曲面
  \[
  3x^2+y^2+z^2=12
  \]
  上点 $M(-1,0,3)$ 处的切平面与平面 $z=0$ 的夹角是（\quad）。

  A. $\dfrac{\pi}{6}$；\quad B. $\dfrac{\pi}{4}$；\quad C. $\dfrac{\pi}{3}$；\quad D. $\dfrac{\pi}{2}$。

  \item 已知 $L$ 是球面
  \[
  x^2+y^2+z^2=1
  \]
  与平面 $x+y+z=0$ 的交线，则第一类曲线积分
  \[
  \int_Lx^2\,ds
  \]
  （\quad）。

  A. $0$；\quad B. $\dfrac{2\pi}{3}$；\quad C. $\dfrac{\pi}{3}$；\quad D. $\dfrac{4\pi}{3}$。

  \item 幂级数
  \[
  \sum_{n=1}^{\infty}\frac{x^n}{n}
  \]
  的和函数是（\quad）。

  A. $s(x)=-\ln(1-x),\ (-1\le x<1)$；\quad
  B. $s(x)=-\ln(x-1),\ (-1<x<1)$；

  C. $s(x)=-\ln(1+x),\ (-1\le x<1)$；\quad
  D. $s(x)=\ln(1-x),\ (-1\le x<1)$。
\end{enumerate}

\subsubsection*{二、填空题（每小题 3 分，共 15 分）}
\begin{enumerate}
  \item 已知
  \[
  f(x,y)=\sin x\ln(y+1)+\cos y\ln(1-x),
  \]
  则 $\left.\dfrac{\partial f}{\partial x}\right|_{(0,0)}=$ \underline{\hspace{3cm}}。
  \item 微分方程
  \[
  y^{(4)}-y=0
  \]
  的通解为 \underline{\hspace{4cm}}。
  \item 设 $D=[0,1;0,2]$，则二重积分
  \[
  \iint_D\sqrt{x+y}\,dx\,dy=
  \]
  \underline{\hspace{3cm}}。
  \item 设 $L$ 是任意一条光滑闭曲线，则
  \[
  \oint_L(2xy+\cos x)\,dx+(x^2+\sin y)\,dy=
  \]
  \underline{\hspace{3cm}}。
  \item 函数项级数
  \[
  \sum_{n=1}^{\infty}\frac1n\left(\frac{x-1}{x}\right)^n
  \]
  的收敛域为 \underline{\hspace{4cm}}。
\end{enumerate}

\subsubsection*{三、计算题（每小题 10 分，共 50 分）}
\begin{enumerate}
  \item 求曲面
  \[
  x^2+y^2+z^2=x
  \]
  的切平面，使它垂直于平面 $x-y-z=2$ 和 $x-y-\dfrac12z=2$。
  \item 求微分方程
  \[
  y''+(4x+e^{2y})(y')^3=0
  \]
  的通解。
  \item 计算三重积分
  \[
  I=\iiint_D(x^2+y^2)\,dx\,dy\,dz,
  \]
  其中 $D$ 是由圆锥面 $x^2+y^2=z^2$ 与平面 $z=1$ 所围成的区域。
  \item 设 $\Sigma$ 为曲面
  \[
  x^2+y^2+z^2=2ax\quad(a>0),
  \]
  计算曲面积分
  \[
  I=\iint_\Sigma(x^2+y^2+z^2)\,dS.
  \]
  \item 将函数
  \[
  f(x)=\frac1{x^2+4x+3}
  \]
  在 $x=1$ 处展成幂级数。
\end{enumerate}

\subsubsection*{四、证明题（本题 10 分）}
设
\[
x=x(y,z),\qquad y=y(x,z),\qquad z=z(x,y)
\]
都是由方程
\[
F(x,y,z)=0
\]
所确定的具有连续偏导数的函数，证明
\[
\frac{\partial x}{\partial y}\frac{\partial y}{\partial z}\frac{\partial z}{\partial x}=-1.
\]

\subsubsection*{五、应用题（本题 10 分）}
设有平面力场
\[
\mathbf F(x,y)=(2xy^3-y^2\cos x)\mathbf i+(1-2y\sin x+3x^2y^2)\mathbf j,
\]
求一质点沿曲线
\[
L:2x=\pi y^2
\]
从点 $O(0,0)$ 运动到点 $A\!\left(\dfrac\pi2,1\right)$ 时变力 $\mathbf F$ 所作的功。

\subsection{答案与解析}

\subsubsection*{一、单项选择题}
\begin{enumerate}
  \item 特征方程 $(r+2)^2=0$，故 $y=(C_1+C_2x)e^{-2x}$。由初值 $C_1=1,\ C_2=5$，选 B。
  \item 偏导连续可推出可微，选 C。
  \item 曲面法向量为 $\nabla F=(6x,2y,2z)$，在 $M$ 处为 $(-6,0,6)$。切平面法向量与 $z=0$ 法向量 $(0,0,1)$ 的夹角余弦为 $\dfrac1{\sqrt2}$，两平面夹角为 $\dfrac\pi4$，选 B。
  \item 交线为单位球上的大圆。由对称性，圆上 $x^2$ 平均值为 $1/3$，弧长为 $2\pi$，故积分为 $\dfrac{2\pi}{3}$，选 B。
  \item 由经典展开
  \[
  -\ln(1-x)=\sum_{n=1}^{\infty}\frac{x^n}{n}
  \]
  且在 $x=-1$ 收敛、$x=1$ 发散，选 A。
\end{enumerate}

\subsubsection*{二、填空题}
\begin{enumerate}
  \item
  \[
  f_x=\cos x\ln(y+1)-\frac{\cos y}{1-x},
  \]
  故 $f_x(0,0)=-1$。
  \item 特征根为 $1,-1,i,-i$，故通解为
  \[
  y=C_1e^x+C_2e^{-x}+C_3\cos x+C_4\sin x.
  \]
  \item 先对 $x$ 积分：
  \[
  \iint_D\sqrt{x+y}\,dx\,dy
  =\frac23\int_0^2\bigl((y+1)^{3/2}-y^{3/2}\bigr)\,dy
  =\frac4{15}(3^{5/2}-2).
  \]
  \item 由 Green 公式，旋度为 $2x-2x=0$，故积分为 $0$。
  \item 令 $t=\dfrac{x-1}{x}=1-\dfrac1x$。级数 $\sum t^n/n$ 的收敛域为 $[-1,1)$，即
  \[
  -1\le1-\frac1x<1.
  \]
  解得
  \[
  x\in\left[\frac12,+\infty\right).
  \]
\end{enumerate}

\subsubsection*{三、计算题}
\begin{enumerate}
  \item 两平面法向量分别为 $\mathbf n_1=(1,-1,-1)$ 与 $\mathbf n_2=(1,-1,-1/2)$。所求切平面同时垂直于它们，其法向量应平行于
  \[
  \mathbf n_1\times\mathbf n_2=(-1,-1,0).
  \]
  曲面法向量为 $(2x-1,2y,2z)$，令其平行于 $(1,1,0)$，得 $z=0,\ 2x-1=2y$。与曲面联立，得两点
  \[
  \left(0,-\frac12,0\right),\qquad \left(\frac45,\frac3{10},0\right).
  \]
  对应切平面为
  \[
  x+y+\frac12=0,\qquad x+y-\frac{11}{10}=0.
  \]
  \item 令 $p=y'$，把 $x$ 看作 $y$ 的函数，利用
  \[
  y''=-\frac{x''}{(x')^3},\qquad y'=\frac1{x'}.
  \]
  原方程化为
  \[
  -x''+4x+e^{2y}=0,
  \]
  即
  \[
  x''-4x=e^{2y}.
  \]
  以 $y$ 为自变量求解得
  \[
  x=C_1e^{2y}+C_2e^{-2y}+\frac14 y e^{2y}.
  \]
  该式为原方程通解的隐式形式。
  \item 用柱坐标，$0\le z\le1,\ 0\le r\le z$。于是
  \[
  I=\int_0^{2\pi}\int_0^1\int_0^z r^2\cdot r\,dr\,dz\,d\theta
  =\frac{\pi}{6}.
  \]
  \item 曲面为以 $(a,0,0)$ 为球心、半径为 $a$ 的球面，且曲面上 $x^2+y^2+z^2=2ax$。由对称性
  \[
  \iint_\Sigma x\,dS=a\cdot4\pi a^2,
  \]
  故
  \[
  I=2a\iint_\Sigma x\,dS=8\pi a^4.
  \]
  \item
  \[
  \frac1{x^2+4x+3}=\frac1{(x+1)(x+3)}
  =\frac12\left(\frac1{x+1}-\frac1{x+3}\right).
  \]
  令 $t=x-1$，则
  \[
  f(x)=\frac12\left(\frac1{t+2}-\frac1{t+4}\right)
  =\sum_{n=0}^{\infty}(-1)^n\left(\frac1{2^{n+2}}-\frac1{2^{2n+3}}\right)t^n,
  \]
  其中 $|t|<2$，即 $|x-1|<2$。
\end{enumerate}

\subsubsection*{四、证明题}
由 $F(x,y,z)=0$ 隐函数求导：
\[
\frac{\partial x}{\partial y}=-\frac{F_y}{F_x},\qquad
\frac{\partial y}{\partial z}=-\frac{F_z}{F_y},\qquad
\frac{\partial z}{\partial x}=-\frac{F_x}{F_z}.
\]
三式相乘立即得
\[
\frac{\partial x}{\partial y}\frac{\partial y}{\partial z}\frac{\partial z}{\partial x}=-1.
\]

\subsubsection*{五、应用题}
记
\[
P=2xy^3-y^2\cos x,\qquad Q=1-2y\sin x+3x^2y^2.
\]
有
\[
P_y=6xy^2-2y\cos x,\qquad Q_x=-2y\cos x+6xy^2,
\]
故力场保守。求势函数：
\[
\Phi_x=P.
\]
对 $x$ 积分得
\[
\Phi=x^2y^3-y^2\sin x+h(y).
\]
再由 $\Phi_y=Q$ 得 $h'(y)=1$，故
\[
\Phi=x^2y^3-y^2\sin x+y.
\]
所作功为
\[
W=\Phi\!\left(\frac\pi2,1\right)-\Phi(0,0)
=\frac{\pi^2}{4}-1+1=\frac{\pi^2}{4}.
\]
